\chapter{Problem Setup and Related Work}
\label{chapter:problem_setup}

\todo{RESTRUCTURE: This chapter should become C3 (Problem Formulation) ONLY. All Related Work content (Sections 2-3) should move to the new standalone C2 Related Work chapter.}

\todo{ADD: Explicit Research Questions R1-R4: (R1) What are the failure modes when agents communicate across trust boundaries? (R2) How does a policy file affect both utility and security? (R3) How should memory be redesigned for multi-party agent interactions? (R4) Can we design an adaptive escalation mechanism that learns social boundaries?}

\todo{MOVE: "Landscape Analysis" section (2.1-2.3) and "The Pulse Niche" section (2.4) to the new C2 Related Work chapter.}

\todo{MOVE: "Key Differentiators" comparison table to C2 Related Work.}

\todo{KEEP: The formal math (agents, files, states, trade-off formulation) -- this is good. Refine the notation.}

\todo{KEEP: Threat model section -- but add citations to actual prompt injection literature (Greshake et al., Perez \& Ribeiro, etc.)}

\section{The Context-Security Trade-off}

The fundamental challenge in networked AI systems can be formalized as the \textbf{Context-Security Trade-off}: achieving high collaborative utility while maintaining absolute security guarantees when agents operate across organizational trust boundaries.

\subsection{Formal Problem Statement}

Let $A$ be a host AI agent owned by user $U$, possessing access to:
\begin{itemize}
    \item $F = \{f_1, f_2, ..., f_n\}$: A set of files (notes, documents)
    \item $S = \{s_1, s_2, ..., s_m\}$: A set of dynamic states (calendar events, emails, todos)
    \item $M_{self}$: Global self-state memory (core facts about $U$)
\end{itemize}

When $A$ interacts with external entity $E_i$ (e.g., an investor, cofounder, or stranger), the agent must:

\begin{enumerate}
    \item \textbf{Maximize Utility}: Complete collaborative tasks $T = \{t_1, t_2, ..., t_k\}$ (e.g., schedule meetings, answer questions, negotiate terms) by leveraging relevant subsets of $F$, $S$, and $M_{self}$.

    \item \textbf{Enforce Security}: Guarantee that $E_i$ can only access authorized data subsets $F' \subseteq F$ and $S' \subseteq S$ despite adversarial prompts attempting to extract $F \setminus F'$ or $S \setminus S'$.

    \item \textbf{Maintain Isolation}: Ensure that memory artifacts from interaction with $E_i$ do not contaminate interactions with $E_j$ for $i \neq j$.
\end{enumerate}

Traditional architectures fail because they conflate these requirements: prompt-based restrictions (``You are not allowed to access file $f_x$'') provide no physical enforcement, while manual permission management for every $(E_i, f_j, s_k)$ triple scales quadratically and reintroduces coordination overhead.

\subsection{Threat Model}

We assume a Byzantine adversary model where external entities $E_i$ may:
\begin{itemize}
    \item Attempt \textbf{prompt injection attacks} to override system instructions
    \item Employ \textbf{social engineering} to manipulate the agent's reasoning
    \item Execute \textbf{data exfiltration} through indirect queries that aggregate protected information
    \item Launch \textbf{confused deputy attacks} by tricking the agent into performing unauthorized actions
\end{itemize}

We do \textit{not} assume compromise of the underlying LLM service, operating system, or network transport layer—these are orthogonal security concerns.

\section{Landscape Analysis}

\subsection{Single-Tenant Agent Frameworks}

Current commercial and academic agent systems prioritize execution capability within isolated environments:

\begin{itemize}
    \item \textbf{OpenClaw} and \textbf{Manus}: High-capability task automation for individual users. These systems excel at workflow orchestration (e.g., ``research this topic and summarize findings'') but assume a single-tenant, trusted execution environment. They provide no architectural primitives for cross-boundary delegation.

    \item \textbf{MemGPT} and \textbf{AgentOS (Wu et al.)}: Advanced memory management systems that maintain long-term context through tiered memory hierarchies (working memory, episodic memory, semantic memory). However, these architectures treat memory as a global resource—there is no concept of relationship-specific isolation or multi-tenant memory partitioning.

    \item \textbf{LangChain/LlamaIndex}: Tooling frameworks for building agents with access to external APIs and data sources. Security is delegated to prompt engineering (``Do not access unauthorized resources''), which fails under adversarial inputs.
\end{itemize}

\textbf{Gap}: These systems are architecturally unsuitable for networked scenarios because they lack physical isolation mechanisms. An agent granted email access for internal use cannot selectively expose only ``Stripe-related emails'' to an external auditor without manual filtering.

\subsection{Agent Communication Protocols}

Academic research has explored Agent-to-Agent (A2A) communication:

\begin{itemize}
    \item \textbf{Multi-Agent Systems (MAS)}: Classical frameworks (JADE, FIPA-ACL) define message-passing protocols for agents to negotiate and collaborate. However, these assume trusted environments where all agents share a common ontology and security model. They do not address zero-trust cross-organizational scenarios.

    \item \textbf{A2A Networks (Wang et al.)}: Recent work on LLM-based agent communication focuses on natural language negotiation. Agents exchange text messages to coordinate tasks (e.g., scheduling meetings). However, this approach is token-inefficient, latency-prone, and vulnerable to hallucination cascades. More critically, it assumes agents have unrestricted access to their owner's data—there is no access governance layer.

    \item \textbf{Federated Learning and Secure Multi-Party Computation}: Privacy-preserving computation techniques allow collaborative learning without data sharing. While conceptually related, these methods are designed for model training, not operational coordination tasks requiring dynamic context sharing.
\end{itemize}

\textbf{Gap}: Existing protocols either assume full trust (MAS) or focus exclusively on message exchange semantics without addressing the underlying access control problem.

\subsection{Data Sharing Platforms}

Commercial file-sharing systems provide secure external access but lack agency:

\begin{itemize}
    \item \textbf{Notion/DocSend/Dropbox}: Secure, permissioned sharing of static documents. Users manually configure read/write access per file. These platforms excel at \textit{static} sharing but offer no dynamic, context-aware interaction. A shared document cannot answer questions, schedule meetings, or execute actions on behalf of its owner.

    \item \textbf{Replika/Character.AI}: Conversational AI companions that maintain long-term memory and personality. These systems demonstrate relationship-aware interaction but operate in purely synthetic domains—they have no connection to real-world operational states (calendars, emails, file systems).
\end{itemize}

\textbf{Gap}: Static sharing platforms provide security without agency; conversational companions provide agency without real-world grounding.

\section{The Pulse Niche: Network-Native AgentOS}

Pulse occupies a unique position in this landscape by synthesizing insights from three domains:

\begin{enumerate}
    \item \textbf{Operating Systems}: The abstraction of ``files'' and ``states'' as mountable resources, capability-based security, and process isolation.

    \item \textbf{Distributed Systems}: Federated network topologies, stateful session management (analogous to TCP), and transitive trust protocols.

    \item \textbf{Agentic AI}: Long-term memory, proactive execution, and natural language interaction.
\end{enumerate}

Unlike single-tenant frameworks, Pulse is \textbf{network-native by design}. Every architectural decision—from memory partitioning to API sandboxing—is driven by the requirement to support secure, stateful interactions across organizational boundaries. This positions Pulse not as an individual productivity tool but as \textbf{foundational infrastructure for an agent economy}—analogous to what TCP/IP provides for computer networks.

\subsection{Key Differentiators}

\begin{table}[h]
\centering
\caption{Comparison of Pulse with existing systems}
\label{table:comparison}
\begin{tabular}{p{3cm} p{2cm} p{2cm} p{2cm} p{2cm}}
\toprule
\textbf{System} & \textbf{Multi-Tenant} & \textbf{Physical Isolation} & \textbf{Real-World State} & \textbf{Network-Native} \\
\midrule
OpenClaw/Manus & No & No & Yes & No \\
MemGPT/AgentOS & No & No & Partial & No \\
A2A Protocols & Yes & No & No & Partial \\
Notion/DocSend & Yes & Yes & No & No \\
Replika/Character.AI & Yes & No & No & No \\
\textbf{Pulse} & \textbf{Yes} & \textbf{Yes} & \textbf{Yes} & \textbf{Yes} \\
\bottomrule
\end{tabular}
\end{table}

\section{Research Questions}

Building on this landscape analysis, we address three core research questions:

\begin{enumerate}
    \item \textbf{RQ1 (Memory)}: How can we architect a memory system that maintains perfect context continuity within each relationship while guaranteeing absolute isolation across relationships?

    \item \textbf{RQ2 (Security)}: Can we design a capability-based execution environment where unauthorized data access is physically impossible, rendering prompt injection attacks ineffective?

    \item \textbf{RQ3 (Usability)}: How can we automatically learn and generalize user privacy preferences across similar relationships to minimize manual permission overhead?
\end{enumerate}

The remainder of this thesis presents our solutions to these questions through the Pulse AgentOS architecture.
